\subsection{Exact local lifting in the arithmetic module quotient}

\begin{statement}{E.42}{Proved}{A tame existence test and finite coordinates for nonzero first inertia}
Use the actual arithmetic quotient of E.38, with
\[
 \mathscr A=\mathbb Z_3[[U]],\quad T=1+U,\quad
 J=(3,U)\mathscr A,\quad H=3\mathbb Z_3\ltimes J.
 \tag{E42.1}
\]
Its law is $(a,v)(b,w)=(a+b,v+T^aw)$, and the based
Galois action is $(a,v(T))\mapsto
(\chi a,\chi v(T^\chi))$. Let $\mathfrak q\nmid3$ be
a prime of $K=\mathbb Q(\zeta)$, put $q=N\mathfrak q$,
and choose local Frobenius $F$ and tame pro-three inertia
$\tau$ with $F\tau F^{-1}=\tau^q$. Here $q\equiv1\pmod3$.
Fix four actual coordinates $a,b,c,d\in3\mathbb Z_3$.
We compute the continuous section cocycles with
\[
 s(\tau)=(a,v),\qquad s(F)=(b,w),\qquad
 v(1)=c,\quad w(1)=d.
 \tag{E42.2}
\]
The notation in (E42.2) records the geometric coordinates
of the section, with the arithmetic generator understood.

For $a=0$ and $c=0$, all solutions have $v=0$, and
$w$ is an arbitrary element of $J$ of augmentation $d$.
For $a=0$ and $c\ne0$, solutions exist exactly when
\[
             \frac{b}{q-1}\in\mathbb Z_3.
 \tag{E42.3}
\]
Every solution then has
$v=cT^{-b/(q-1)}$, and again $w$ is arbitrary with the
specified augmentation. For any negative valuation of
$b/(q-1)$, the obstruction holds for every nonzero $c$,
regardless of its three-adic divisibility.

For $a\ne0$, set $m=3^{v_3(a)}$ and list the orbits
$O_1,\ldots,O_e$ of the actual affine permutation
\[
 \mathbb Z/m\longrightarrow\mathbb Z/m,
                  \quad k\longmapsto qk+b.
 \tag{E42.4}
\]
Use representatives $0\le k<m$, and define
\[
 \begin{split}
 P_i(T)&=\sum_{k\in O_i}T^k,\quad
 \ell_i=|O_i|,\quad \sigma_i=\sum_{k\in O_i}k,\\
 \delta_i&=\frac{(q-1)\sigma_i+
                      (b-a(q-1)/2)\ell_i}{a}\in\mathbb Z_3.
 \end{split}
 \tag{E42.5}
\]
The displayed integrality is part of the assertion.
Solutions exist exactly when the following finite
two-row integral system has a solution
$(z_1,\ldots,z_e,h_0)\in\mathbb Z_3^{e+1}$:
\[
 \sum_i\ell_i z_i=c,\qquad
 \sum_i\delta_i z_i+(q-1)h_0=d.
 \tag{E42.6}
\]
This is an explicit Smith test over $\mathbb Z_3$.
For every solution, put
\[
 \begin{split}
 n&=T^a-1,\qquad S=\sum_{j=0}^{q-1}T^{aj},\\
 P&=\sum_i z_iP_i,\qquad h=h_0+U h_1,\quad h_1\in\mathscr A,\\
 D(P)&=\frac{qT^bP(T^q)-SP(T)}{nS}\in\mathscr A,\\
 v&=P+nh,\qquad
 w=D(P)+qT^bh(T^q)-h(T).
 \end{split}
 \tag{E42.7}
\]
These are all the actual section cocycles in (E42.2).
Modulo geometric conjugations preserving all four fixed
coordinates, their classes are in bijection with the
solutions of (E42.6). Every nonempty such set is a
torsor for the free module $\mathbb Z_3^{e-1}$, with
coordinates supplied by that finite matrix.

For $a=0$ the conjugations preserving the four coordinates
are also explicit: they leave $v$ fixed and send
$w$ to $w+u-qT^bu(T^q)$, with $u\in U\mathscr A$.
Thus the formulas specify both the actual cocycles and
their equivalence relation in this case as well.

\emph{Manuscript: \S51.16, folios 733v--736v, scans 172--178;
Part E(c),(d) and the fourth exactness follow-up. The
arithmetic quotient maps $N/N''\to H$ and
$N/\gamma_3I_{\rm marked}\to H$, with their kernels,
are precisely (E38.5)--(E38.9).}
\end{statement}
\begin{proof}
The action on $H$ is the generator computation of E.38.
Over $K_{\mathfrak q}$ its cyclotomic image is pro-three.
Consequently a section cocycle factors through the maximal
pro-three quotient of the local Galois group: its graph
maps into the pro-three group formed by $H$ and that
cyclotomic image. Killing wild inertia, which is pro-$p$
for $p\ne3$, and the prime-to-three tame factors gives
the actual presentation
\[
 \langle F,\tau\mid F\tau F^{-1}=\tau^q\rangle^{\mathrm{pro}\text{-}3}.
 \tag{E42.8}
\]
The unramified factor and the remaining inertia are each
$\mathbb Z_3$, with conjugation multiplication by $q$.
Conversely images of these two generators satisfying
the relation define a continuous homomorphism by this
presentation. Thus the equation below tests actual
continuous local sections.

Compute the geometric coordinates of $F\tau$ and
$\tau^qF$ in $H$ semidirect the arithmetic action.
Their first coordinates are both $b+qa$. Their second
coordinates give exactly
\[
 (T^{qa}-1)w(T)=qT^bv(T^q)-S_q(T^a)v(T),
 \quad S_q(T^a)=\sum_{j=0}^{q-1}T^{aj}.
 \tag{E42.9}
\]
The scalar $q$ in the first term and the translation
$T^b$ both come from the genuine action.

For $a=0$ this reduces to $T^bv(T^q)=v(T)$.
Work first in $\mathbb Q_3[[U]]$ and substitute
$T=\exp z$. Multiplication by
$\exp(bz/(q-1))$ converts the equation to invariance
under $z\mapsto qz$. Since $q^j-1\ne0$ for every
$j>0$, the resulting series is constant. Hence
$v=cT^{-b/(q-1)}$, including $v=0$ for $c=0$.
For an exponent $r\in\mathbb Z_3$, the coefficients
$\binom rj$ are integral, by continuity from integer
exponents. For $v_3(r)=-s<0$, every factor $r-i$,
$i\in\mathbb Z$, has that same valuation. Therefore
\[
 v_3\!\left(c\binom rj\right)
             =v_3(c)-sj-v_3(j!),
 \tag{E42.10}
\]
which is negative for large $j$ whenever $c\ne0$.
This proves the exact integral alternative (E42.3).
The term containing $w$ vanishes in (E42.9), so
there is no further restriction on $w$.

Now let $a=mu$ with $u\in\mathbb Z_3^\times$. The
quotient
$((T^m)^u-1)/(T^m-1)$ is an element of $\mathscr A$
with constant term $u$, hence is a unit. The ideals
$(n)$ and $(T^m-1)$ are consequently equal. The
actual quotient map $T\mapsto\overline T$ gives
\[
 \mathscr A/(n)\xrightarrow{\sim}
              \mathbb Z_3[\mathbb Z/m].
 \tag{E42.11}
\]
For example, division by the monic distinguished
polynomial $(1+U)^m-1$ gives the unique representatives
of degree less than $m$; this also proves the claimed
topological isomorphism. The factor $S_q(T^a)$ has
augmentation $q\in\mathbb Z_3^\times$, so it too
is a unit in $\mathscr A$.

Modulo $n$, equation (E42.9) becomes
$T^bv(T^q)=v(T)$. In the displayed group-ring basis
this is precisely invariance under the permutation
(E42.4). A vector is invariant exactly when its
coefficients are constant along each orbit. Thus
its invariant submodule has the free basis $P_i$,
and every permissible $v$ has a unique expression
$v=P+nh$ as in (E42.7). Invariance makes the
numerator defining $D(P)$ divisible by $n$; division
by the unit $S$ keeps it integral. Substituting
$v=P+nh$ in (E42.9) and using
$n(T^q)=n(T)S(T)$ gives the exact formula for $w$
in (E42.7). Since $\mathscr A$ is an integral domain,
this $w$ is unique for the chosen $v$.

Differentiate the numerator and denominator of
$D(P_i)$ at $T=1$. The numerator has value zero,
the derivative of $nS$ is $qa$, and the numerator
derivative is
\[
 q\bigl((q-1)\sigma_i+
                  (b-a(q-1)/2)\ell_i\bigr).
\]
It follows that $D(P_i)(1)=\delta_i$.
The integrality of $D(P_i)$ proves that of this
formula even when $a$ is divisible by a large
power of three. Taking augmentations in (E42.7)
now gives exactly (E42.6). The required values
$c,d\in3\mathbb Z_3$ ensure $v,w\in J$, so these
are sections of $H$ itself.

We check the conjugacy assertion without changing
the specified coordinates. Conjugation by an element
$(s,u)\in H$ changes the first Frobenius coordinate
$b$ by $(1-q)s$. Equality with the fixed $b$ forces
$s=0$. With $s=0$, the inertia and Frobenius module
coordinates change respectively by
\[
 v\longmapsto v-nu,\qquad
 w\longmapsto w+u-qT^bu(T^q).
 \tag{E42.12}
\]
Preserving $d$ forces $(1-q)u(1)=0$, hence
$u\in U\mathscr A$. For $a\ne0$, (E42.12)
changes $h$ to $h-u$, fixes every $z_i$ and $h_0$,
and removes the arbitrary $Uh_1$. These are all
conjugations preserving the four coordinates.
The matrix in (E42.6) has rank two: its last column
is $(0,q-1)^{\mathsf T}$ and each $\ell_i$ is
nonzero. Its kernel is therefore a free
$\mathbb Z_3$-module of rank $e-1$. This proves
the full assertion about those classes. For $a=0$
the same computation gives the equivalence relation
stated after (E42.7). Its linear image is closed,
since $U\mathscr A$ is compact and the map is
continuous.

The program \texttt{e\_tame\_module.py} checks
$96$ finite orbit systems, with $q=4,7,19,37$,
$m=3,9,27$, and both units $1,2$ multiplying $m$.
It verifies the two-row Smith valuations, both
augmentations, and the original relation through
degree $18$ by exact rational coefficients integral
at three. It also checks $28$ zero-first-inertia
cases, including the forced coefficient $1729/384$
of E.38. The all-degree conclusions follow from
the ideal, permutation and power-series arguments
above; no finite precision is used to infer them.
\end{proof}

